\section{A linear map is determined by its building blocks}\label{ch10:sec:linearmaps}
Section~\ref{ch06:sec:structure} said that a map preserves addition when it sends the sum of two inputs to the sum of their outputs. A map $T:\RR^n\to\RR^m$ is \emph{linear} when it preserves both addition and multiplication by scalars:
\[
T(x+y)=T(x)+T(y)\qquad\text{and}\qquad T(cx)=cT(x)
\]
for all vectors $x,y\in\RR^n$ and all scalars $c$. The two conditions can be combined into one: $T$ is linear exactly when
\[
T(ax+by)=aT(x)+bT(y)\qquad\text{for all $x,y\in\RR^n$ and all $a,b\in\RR$.}
\]
Indeed, if $T$ satisfies the two conditions, then $T(ax+by)=T(ax)+T(by)=aT(x)+bT(y)$. Conversely, the combined condition with $a=b=1$ is preservation of addition, and with $b=0$ it says $T(ax)=aT(x)$, because $ax+0y=ax$ and $aT(x)+0T(y)=aT(x)$.

Taking $c=0$ in the second condition shows that a linear map sends the zero vector to the zero vector: for any vector $x$, we have $0x=0$, so $T(0)=T(0x)=0\,T(x)=0$.

A linear map sends every linear combination to the same combination of the outputs. For three vectors, first split off the last term using preservation of addition. Then apply the combined condition to the first piece and preservation of scaling to the second:
\[
T(a_1v_1+a_2v_2+a_3v_3)=T(a_1v_1+a_2v_2)+T(a_3v_3)=a_1T(v_1)+a_2T(v_2)+a_3T(v_3).
\]
Splitting off one term at a time in the same way works for any number $k$ of terms (formally, this is an induction on $k$, as in Chapter~\ref{ch:induction}):
\[
T(a_1v_1+\cdots+a_kv_k)=a_1T(v_1)+\cdots+a_kT(v_k).
\]

Now apply this to the standard basis. For a linear map $T:\RR^2\to\RR^m$, writing $(x_1,x_2)=x_1e_1+x_2e_2$ gives
\[
T(x_1,x_2)=T(x_1e_1+x_2e_2)=x_1T(e_1)+x_2T(e_2).
\]
Here $T(x_1,x_2)$ is short for $T\bigl((x_1,x_2)\bigr)$, the output at the vector $(x_1,x_2)$; we leave out one pair of brackets. In the same way, writing $x=x_1e_1+\cdots+x_ne_n$ shows that every linear map $T:\RR^n\to\RR^m$ satisfies
\[
T(x)=x_1T(e_1)+x_2T(e_2)+\cdots+x_nT(e_n).
\]
So the $n$ vectors $T(e_1),\ldots,T(e_n)$ determine the output $T(x)$ at every one of the infinitely many inputs $x$. This is why a linear map can be recorded in a finite table of numbers.

Suppose, for example, that $T:\RR^2\to\RR^2$ is linear, with $T(e_1)=(1,1)$ and $T(e_2)=(1,-1)$. Then
\[
T(x_1,x_2)=x_1(1,1)+x_2(1,-1)=(x_1+x_2,\,x_1-x_2).
\]
The two outputs $T(e_1)$ and $T(e_2)$ are the basis vectors $w_1$ and $w_2$ of Section~\ref{ch10:sec:vectors}. A \emph{matrix} is a rectangular table of numbers. We record $T$ by the matrix whose first column is $T(e_1)$ and whose second column is $T(e_2)$:
\[
A=\begin{pmatrix}1&1\\1&-1\end{pmatrix},\qquad
Ax=\begin{pmatrix}x_1+x_2\\x_1-x_2\end{pmatrix}.
\]
Here $Ax$ stands for $T(x)$, the output of the map that $A$ records. We sometimes write a vector horizontally, as $(x_1,x_2)$, and sometimes vertically, as a column, as in this display. Both name the same ordered list. When we apply a matrix to a vector, as in $Ax$, we write the vector as a column to the right of the matrix, because the row rule below pairs each row of the matrix with that column.

An \emph{$m$-by-$n$ matrix} has $m$ rows and $n$ columns. We write $A_{ij}$ for the entry in row $i$ and column $j$; the first subscript always names the row. For the matrix $A$ above, $A_{12}=1$ and $A_{22}=-1$. In general, an $m$-by-$n$ matrix $A$ acts on a vector $x\in\RR^n$ and produces the vector $Ax\in\RR^m$ whose $i$th coordinate is
\[
(Ax)_i=\sum_{j=1}^n A_{ij}x_j=A_{i1}x_1+A_{i2}x_2+\cdots+A_{in}x_n\qquad(1\le i\le m).
\]
In words: to find the $i$th coordinate of $Ax$, multiply the entries of row $i$ by the matching coordinates of $x$, and add the products. We call this the \emph{row rule}. Section~\ref{ch10:sec:inner} gives this kind of sum of products its usual name, the dot product.

For example, take
\[
M=\begin{pmatrix}2&-1&3\\1&0&1\end{pmatrix}\qquad\text{and}\qquad x=(4,5,0).
\]
The first row gives $2\cdot4+(-1)\cdot5+3\cdot0=3$, and the second row gives $1\cdot4+0\cdot5+1\cdot0=4$, so $Mx=(3,4)$. The input must have three coordinates, because each row has three entries; the output has two coordinates, one for each row. Here $m=2$ and $n=3$. In general the numbers $m$ and $n$ may differ, and a matrix with $m=n$ is called \emph{square}.

The row rule agrees with the way $A$ was built from its columns. For any $m$-by-$n$ matrix $A$, the $i$th coordinate of the combination
\[
x_1(\text{column $1$ of $A$})+x_2(\text{column $2$ of $A$})+\cdots+x_n(\text{column $n$ of $A$})
\]
is $x_1A_{i1}+x_2A_{i2}+\cdots+x_nA_{in}$, which is exactly $(Ax)_i$. So $Ax$ is this combination of the columns, with the coordinates of $x$ as coefficients. In the example,
\[
Mx=4(2,1)+5(-1,0)+0(3,1)=(8,4)+(-5,0)+(0,0)=(3,4),
\]
as before. Taking $x=e_j$ shows that $Ae_j$ is column $j$ of $A$. This is why we put $T(e_j)$ in column $j$. For any linear map $T:\RR^n\to\RR^m$, build the $m$-by-$n$ matrix $A$ whose column $j$ is $T(e_j)$, for $j=1,\ldots,n$. Since $Ax$ is the combination of the columns with coefficients $x_1,\ldots,x_n$, we get $Ax=x_1T(e_1)+\cdots+x_nT(e_n)$, and this equals $T(x)$. So the matrix $A$ records $T$ exactly; it is called the \emph{matrix of} $T$.

Conversely, every $m$-by-$n$ matrix $A$ gives a linear map $x\mapsto Ax$ from $\RR^n$ to $\RR^m$. Take vectors $x,y\in\RR^n$ and scalars $a,b$. The $j$th coordinate of $ax+by$ is $ax_j+by_j$, so by the row rule the $i$th coordinate of $A(ax+by)$ is
\[
\sum_{j=1}^n A_{ij}(ax_j+by_j)=a\sum_{j=1}^n A_{ij}x_j+b\sum_{j=1}^n A_{ij}y_j,
\]
where we multiplied out each bracket and collected the terms containing $a$ and those containing $b$. The right side is the $i$th coordinate of $aAx+bAy$. This holds for every $i$, so $A(ax+by)=aAx+bAy$.

The two constructions undo each other. Starting from a matrix $A$, the matrix of the map $x\mapsto Ax$ has column $j$ equal to $Ae_j$, which is column $j$ of $A$; so we get $A$ back. Starting from a linear map $T$, its matrix $A$ satisfies $Ax=T(x)$ for every $x$, so we get $T$ back, and no other matrix does this: if $Bx=T(x)$ for every $x$, then column $j$ of $B$ is $Be_j=T(e_j)$, so $B=A$. Thus each linear map from $\RR^n$ to $\RR^m$ has exactly one matrix, each $m$-by-$n$ matrix comes from exactly one linear map, and the two are just different descriptions of the same object. For this reason we often say that a matrix maps $\RR^n$ to $\RR^m$.

\pauseq{Is $F(x_1,x_2)=(x_1+1,x_2)$ linear?}{No. A linear map sends the zero vector to the zero vector, but $F(0,0)=(1,0)$. The map $F$ is a \emph{translation}: it adds the fixed vector $(1,0)$ to every input. It does preserve differences, since $F(x)-F(y)=x-y$, but it moves the zero vector, and that alone rules out linearity.}
